\documentclass[11pt]{article}
\usepackage[margin=0.95in]{geometry}
\usepackage{amsmath,amssymb}
\usepackage{booktabs}
\usepackage{slashed}
\usepackage[colorlinks=true,linkcolor=blue,urlcolor=blue]{hyperref}
\usepackage{parskip}
\setlength{\parindent}{0pt}

\title{\vspace{-2em}Comments on\\[0.3em]
\emph{Fully accounted NLO electromagnetic radiative corrections\\
to the width of $J/\psi$ decay}}
\author{V. Kubarovsky}
\date{17 September 2026}

\begin{document}
\maketitle
\vspace{-1em}

\noindent
Equation, table and reference numbers refer to the current draft. ``The note''
is \emph{Radiative Corrections to the $J/\psi$ photoproduction reaction};
Ref.~[10] of the draft is Ehlotzky \& Mitter, Nuovo Cimento \textbf{55A} (1968)
181, abbreviated E\&M below.

Sections~1--4 are small textual points. Sections~5--8 are substantive.

\section{Eq.~(1) and Eq.~(17): replace $E_V$ by $m_V$}

Both read $1/(2E_V)$. The calculation is done in the $J/\psi$ rest frame, where
$E_V=m_V$ identically, and the text already says so. Writing $m_V$ removes a
symbol that never takes any other value and makes the normalisation of Eq.~(1)
directly comparable with the standard two-body form.

\section{Eq.~(6): use the factorised form, and add the ultrarelativistic limit}

Eq.~(6) is correct --- algebraically identical to Eq.~(7) of the note (verified
symbolically; the difference is exactly zero). But the factorised form is easier
to read and to check dimensionally, since the mass enters only through
dimensionless ratios:
\begin{equation*}
\Gamma_B=\frac{4\pi\alpha^{2}}{f_V^{2}}\,\frac{m_V}{3}
\left(1+\frac{2m_\ell^{2}}{m_V^{2}}\right)\sqrt{1-\frac{4m_\ell^{2}}{m_V^{2}}}\,.
\end{equation*}
Please also add the ultrarelativistic limit, Eq.~(8) of the note, which is the
form actually used to fix $f_V$ from the measured leptonic width:
\begin{equation*}
\Gamma_B=\frac{4\pi\alpha^{2}}{f_V^{2}}\,\frac{m_V}{3}\,.
\end{equation*}
Worth stating explicitly: the correction connecting the two is \emph{second} order
in the mass ratio. With $x=m_\ell^{2}/m_V^{2}$,
\begin{equation*}
(1+2x)\sqrt{1-4x}=1-6x^{2}+O(x^{3}),
\end{equation*}
so the linear term cancels. This is why the massless limit is so accurate:
$4\times10^{-15}$ for $J/\psi\to e^+e^-$, $8\times10^{-6}$ for
$J/\psi\to\mu^+\mu^-$. It reaches $2\times10^{-3}$ only for $\rho\to\mu^+\mu^-$.

\section{Eq.~(4): the polarization sum}

As printed,
$\sum_{\rm pol}\varepsilon_\alpha(p_V)\varepsilon_\beta(p_V)=-g_{\alpha\beta}$
is the polarization sum for a \emph{massless} vector. For a massive vector meson
it is $-g_{\alpha\beta}+p_\alpha p_\beta/m_V^{2}$. The extra term does drop here,
because it contracts the lepton current with $p_V=p_1+p_2$ and
$\bar u(p_1)(\slashed p_1+\slashed p_2)v(p_2)=0$ by the Dirac equations --- so
Eq.~(5) and everything after it are correct. But the identity as written is not
true on its own, and the surrounding text calls it ``the photon'' while the sum
is over the three $J/\psi$ polarizations. A clause such as ``the
$p_\alpha p_\beta$ term is dropped by current conservation'' would fix it.

\section{Reference [9]}

\texttt{[9] P. Chatagnon, e-Print: 2602.22128 [hep-ex] (2026)} is now published.
Please replace with
\texttt{[9] P. Chatagnon, Phys. Rev. C \textbf{113} (2026) 6, 065203}.

\section{Please add the closed-form total correction, for both channels}

The paper gives $\delta_{self}$~(8), $\delta_{vert}$~(12), $\delta_{soft}$~(13),
and $\delta_{hard}$ only as an unevaluated integral~(14), then the exact
result~(28). There is no compact closed form for the total correction --- the
quantity an experiment applies, and what a reader comes to this paper looking
for. These are Eqs.~(19) and~(20) of the note, i.e.\ Eq.~(26) of Ref.~[10] with
the missing term restored.

\medskip
\textbf{Electrons}, $J/\psi\to e^+e^-$:
\begin{equation*}
\boxed{\;
\delta(m_V,k_{max})=-\frac{4\alpha}{\pi}
\left[
\left(\frac12-\ln\frac{m_V}{m_e}\right)\left(\ln r_{max}+\frac13\right)
+\frac12\left(L_2(r_{max})-\frac{\pi^{2}}{6}\right)
+\frac{17}{72}+R(r_{max})
+\frac{r_{max}^{2}}{16}-\frac{3}{8}r_{max}
\right]\;}
\end{equation*}

\textbf{Muons}, $J/\psi\to\mu^+\mu^-$:
\begin{equation*}
\boxed{\;
\begin{aligned}
\delta(m_V,k_{max})=-\frac{4\alpha}{\pi}
\Bigg[&
\left(\frac12-\ln\frac{m_V}{m_\mu}\right)\left(\ln r_{max}+\frac13\right)
+\frac12\left(L_2(r_{max})-\frac{\pi^{2}}{6}\right)
+\frac{17}{72}\\[2pt]
&+\left(\frac{5}{18}-\frac13\ln\frac{m_V}{m_e}\right)
+R(r_{max})
+\frac{r_{max}^{2}}{16}-\frac{3}{8}r_{max}
\Bigg]
\end{aligned}
\;}
\end{equation*}
with, in each case,
\begin{align*}
r_{max}&=\frac{2k_{max}}{m_V}, &
R(r_{max})&=\left(r_{max}-\frac{r_{max}^{2}}{4}-\frac34\right)
\left[\ln\!\left(\frac{m_V}{m_\ell}\sqrt{1-r_{max}}\right)-\frac12\right], &
L_2(x)&=-\int_0^x\frac{\ln(1-t)}{t}dt,
\end{align*}
$m_\ell$ being the lepton of the channel concerned. The extra term in the muon
expression, $\frac{5}{18}-\frac13\ln\frac{m_V}{m_e}$, is the electron-loop vacuum
polarization; it is exactly $-\frac{\pi}{4\alpha}$ times $\delta_{self}$ of
Eq.~(8) evaluated with $m_e$, and it correctly keeps $m_e$, not $m_\mu$. Here
$r_{max}$ is E\&M's variable, $k_{max}/E$ with $E=m_V/2$; it is identical to the
$r_{cut}=\omega_{cut}/E_1$ already defined in this paper. The last two terms of
each bracket are the ones absent from [10].

\subsection*{Where it belongs}

Section~IV, immediately after Eq.~(14) --- the closed form \emph{is} the
evaluation of that integral combined with Eqs.~(8), (12) and~(13).
\begin{itemize}\itemsep2pt
\item Section~IV opens by promising ``compact expressions for all contributions
      to decay width: both for virtual and real bremstrahlung ones.'' Compact
      expressions follow for every individual contribution, but not for the sum,
      which is the only physical combination.
\item Eq.~(16), $\delta_{FSIC}=3\alpha/4\pi$, then becomes an immediate corollary:
      it is the $r_{max}\to1$ limit of the closed form, less $\delta_{self}$.
\item Figure~4 plots exactly this quantity. At present the reader is shown the
      curve without being given the function.
\end{itemize}
Section~IV also states that the URA is ``applicable for $\ell=e$'' and that the
mass-dependent treatment is ``crucial for \ldots the muon channel''. We would
therefore present the electron expression as the working formula of Section~IV,
and the muon one alongside with the $(m_\mu/E)^{2}$ error flagged and Eq.~(28)
named as the fallback.

\subsection*{Status of the approximation}

These expressions are \emph{not} exact in the lepton mass, and not the
$m_\ell=0$ limit either: the mass is kept in the logarithms, while terms
suppressed by $(m_\ell/E)^{2}$, $E=m_V/2$, are dropped. Ref.~[10] states this
just above its Eq.~(22) (``\ldots neglecting terms of order $(m/E)^{2}$'').
\begin{center}
\begin{tabular}{lccc}
\toprule
 & $(m_\ell/E)^{2}$ & $\delta_{self}$: Eq.~(8)$-$Eq.~(9) & total at $r_{max}=1$: this form $-$ Eq.~(28)\\
\midrule
$e^+e^-$     & $1.1\times10^{-7}$ & $+3\times10^{-10}$  & $+1\times10^{-7}$\\
$\mu^+\mu^-$ & $4.7\times10^{-3}$ & $+1.1\times10^{-5}$ & $-7\times10^{-5}$\\
\bottomrule
\end{tabular}
\end{center}
So for electrons the closed form is exact at any precision an experiment will
reach; for muons it is good to $\sim10^{-4}$ absolute. Presenting it alongside
Eq.~(28) lets a reader use the simple expression where it suffices and the exact
one where it does not.

Checks these expressions pass: they reproduce Table~I of the draft to all six
digits at every $r_{cut}$, in both channels, and are independent of $\omega$ by
construction; at $r_{max}=1$ they give exactly $\delta_{self}+3\alpha/4\pi$; and
for $J/\psi\to e^+e^-$ at $r_{max}=1$ the value is $0.026135$, matching the
Bardin--Shumeiko result of Table~III.

\section{The missing term originates in Ref.~[10]}

E\&M Eq.~(23), the differential spectrum, is correct and agrees with Eq.~(14) of
the draft. The term was lost when that spectrum was integrated analytically in
1968, so Eqs.~(24) and~(26) of [10] are both incomplete. The note transcribed
those equations faithfully and inherited the omission; it has since been
corrected.

\textbf{Where to record this:} in Section~IV, immediately after the closed-form
expressions of Section~5 above --- right where the new formulas will appear,
following Eq.~(14). A single paragraph there attributes the formula, records the
correction, and explains the numerical comparison with [10] all at once:

\begin{quote}
The expressions above are Eq.~(26) of Ref.~[10] with the term
$(\alpha/\pi)(3r_{max}/2-r_{max}^{2}/4)$ restored. The differential spectrum,
Eq.~(23) of [10], is correct and agrees with our Eq.~(14); the term was lost in
its analytic integration, so that Eq.~(24), Eq.~(26) and Table~I of [10] are all
low by that amount --- $5\alpha/4\pi=0.0029$ at $r_{max}=1$, and below the
precision quoted in [10] for $r_{max}\lesssim0.1$.
\end{quote}

\section{Table~I of Ref.~[10] was computed with the incomplete formula}

This matters for the validation in Section~IV, because Table~I of [10] is not a
correct benchmark: it was computed from E\&M's own Eq.~(26) and carries the same
omission. The evidence is not one line but the whole table --- the $\delta$
column at $r_{max}=1.00$, all fourteen mass rows, both channels, 28 entries:

\begin{center}
\begin{tabular}{lcc}
\toprule
agreement with the $\delta$ printed in Table~I of [10] & electrons & muons\\
\midrule
E\&M Eq.~(26) \emph{as published} (term absent) & 12/14 rows exact & 12/14 rows exact\\
E\&M Eq.~(26) \emph{corrected}                  & 0/14             & 0/14\\
\bottomrule
\end{tabular}
\end{center}

The two misses in each channel are last-digit rounding, not disagreement; the
corrected formula misses every row by $\simeq0.003$. The agreement reported in
the draft at $\omega/m_e=0.01$ therefore holds only in the small-$r_{cut}$ rows,
where the missing term falls below the three decimals quoted in [10].

Table~I of the draft is itself correct --- it follows from Eqs.~(8), (12), (13),
(14), which are all sound. It is only the benchmark that is at fault. Two ways to
handle this, not exclusive:

\textbf{(1) State it}, using the paragraph of Section~6 above. The sentence
claiming the comparison ``succeeds'' should become a pointer to it, since as
written it asserts agreement with a table that is in error.

\textbf{(2) Recompute Table~I at $m_V=M(J/\psi)$.} Since [10] can no longer serve
as the validation target, a table at the physical mass is more useful: it
demonstrates the $\omega$-independence just as well, and the numbers are ones a
$J/\psi$ experiment can use. For $m_V=3.0969$~GeV, in the format of the draft:

\begin{center}
\begin{tabular}{cc rrr c rrr}
\toprule
& & \multicolumn{3}{c}{$e^-e^+$} & & \multicolumn{3}{c}{$\mu^-\mu^+$}\\
\cmidrule(lr){3-5}\cmidrule(lr){7-9}
$r_{cut}$ & $\lg(\omega/m_\ell)$ & $\delta_R$ & $\delta_H$ & $\delta$ & &
$\delta_R$ & $\delta_H$ & $\delta$\\
\midrule
      & $-6$ & $-1.577194$ & $1.481750$ & $-0.095444$ & & $-0.382419$ & $0.376596$ & $-0.005823$\\
0.10  & $-4$ & $-1.225925$ & $1.130481$ & $-0.095444$ & & $-0.259274$ & $0.253452$ & $-0.005823$\\
      & $-2$ & $-0.874655$ & $0.779212$ & $-0.095443$ & & $-0.136130$ & $0.130328$ & $-0.005801$\\
\midrule
      & $-6$ & $-1.577194$ & $1.540558$ & $-0.036636$ & & $-0.382419$ & $0.396794$ & $+0.014375$\\
0.25  & $-4$ & $-1.225925$ & $1.189288$ & $-0.036636$ & & $-0.259274$ & $0.273649$ & $+0.014375$\\
      & $-2$ & $-0.874655$ & $0.838019$ & $-0.036636$ & & $-0.136130$ & $0.150526$ & $+0.014396$\\
\midrule
      & $-6$ & $-1.577194$ & $1.576919$ & $-0.000275$ & & $-0.382419$ & $0.408881$ & $+0.026462$\\
0.50  & $-4$ & $-1.225925$ & $1.225649$ & $-0.000275$ & & $-0.259274$ & $0.285736$ & $+0.026462$\\
      & $-2$ & $-0.874655$ & $0.874380$ & $-0.000275$ & & $-0.136130$ & $0.162613$ & $+0.026483$\\
\midrule
      & $-6$ & $-1.577194$ & $1.603329$ & $+0.026135$ & & $-0.382419$ & $0.416436$ & $+0.034017$\\
1.00  & $-4$ & $-1.225925$ & $1.252060$ & $+0.026135$ & & $-0.259274$ & $0.293291$ & $+0.034017$\\
      & $-2$ & $-0.874655$ & $0.900791$ & $+0.026136$ & & $-0.136130$ & $0.170168$ & $+0.034038$\\
\bottomrule
\end{tabular}
\end{center}
$\delta_R$ for the muon channel includes both self-energy loops. The
$r_{cut}=1.00$ electron entry, $0.026135$, equals $\delta_{self}+3\alpha/4\pi$
and matches Table~III.

\section{The upper limit of $r_{cut}$ is not 1}

The text before Eq.~(16) says the cut energy takes its maximum ``at (15) or
$r_{cut}=r_{max}=1$'', and Table~III is captioned $(r_{max}=1)$. That is
inconsistent with Eq.~(15). With $r_{cut}=2\omega_{cut}/m_V$ and
$\omega_{max}=m_V/2-2m_\ell^{2}/m_V$,
\begin{equation*}
r_{cut}^{\;\max}=\frac{2\omega_{max}}{m_V}=1-\frac{4m_\ell^{2}}{m_V^{2}}=\beta^{2},
\end{equation*}
which is $0.99535$ for muons, not 1. ($1-1.1\times10^{-7}$ for electrons.)

This is not only notational. Eq.~(29) defines
$\rho_1(\omega_{cut})=\sqrt{m_V(m_V-2\omega_{cut})}+\sqrt{m_V(m_V-2\omega_{cut})-4m_\ell^{2}}$.
Setting $\omega_{cut}=m_V/2$, i.e.\ literally $r_{cut}=1$, makes the second
radicand $-4m_\ell^{2}<0$ and $\rho_1$ imaginary. At the true endpoint the
radicand vanishes exactly and $\rho_1(\omega_{max})=2m_\ell$, the correct finite
limit; Eq.~(28) is well defined up to $\omega_{max}$ and no further.

We take it that Table~III was evaluated at $\omega_{cut}=\omega_{max}$ and only
the label is wrong. We suggest replacing ``$r_{cut}=r_{max}=1$'' and the
Table~III caption by ``$\omega_{cut}=\omega_{max}$'' or ``$r_{cut}=\beta^{2}$''.
The axis of Fig.~4 is already correctly drawn as $\omega_{cut}/\omega_{max}$.

\section{Status of the note}

The note has been corrected: the missing term in its Eqs.~(14), (15), (19), (20);
the leading logarithm of its Eq.~(20), which had $m_e$ where $m_\mu$ belongs; the
phase-space factor of its Eq.~(6); and the maximum photon energy, which now
carries the exact endpoint and the bound $r_{max}\le\beta^{2}$. Its Fig.~8 has
been redrawn against $r_{max}$ over the full physical range, so it now coincides
with Fig.~4 of this draft --- the two abscissae are related by
$r_{max}=(\omega_{cut}/\omega_{max})\beta^{2}$, and the curves agree to machine
precision in both channels. A paragraph giving the correspondence between the
variables of the two papers ($\varepsilon\leftrightarrow\omega$,
$k_{max}\leftrightarrow\omega_{cut}$, $r_{max}\leftrightarrow r_{cut}$,
$\ln(m_V/m_\ell)=L/2$) has been added.

\end{document}
